\documentclass[12pt, a4paper]{article}

% ── Packages ──────────────────────────────────────────────────────────────────
\usepackage[utf8]{inputenc}
\usepackage[T1]{fontenc}
\usepackage[english]{babel}
\usepackage{lmodern}
\usepackage{geometry}
\geometry{top=2.5cm, bottom=2.5cm, left=2.8cm, right=2.5cm}
\usepackage{graphicx}
\usepackage{amsmath}
\usepackage{amssymb}
\usepackage{booktabs}
\usepackage{tabularx}
\usepackage{multirow}
\usepackage{longtable}
\usepackage{hyperref}
\usepackage{xcolor}
\usepackage{listings}
\usepackage{caption}
\usepackage{subcaption}
\usepackage{float}
\usepackage{enumitem}
\usepackage{algorithm}
\usepackage{algpseudocode}
\usepackage{setspace}
\usepackage{titlesec}
\usepackage{fancyhdr}
\usepackage[numbers,sort&compress]{natbib}
\usepackage{url}
\usepackage{parskip}
\usepackage{array}
\usepackage{diagbox}

% ── Hyperref setup ─────────────────────────────────────────────────────────────
\hypersetup{
    colorlinks=true,
    linkcolor=blue!70!black,
    citecolor=green!40!black,
    urlcolor=blue!60!black,
}

% ── Listings ──────────────────────────────────────────────────────────────────
\lstset{
    basicstyle=\ttfamily\small,
    breaklines=true,
    frame=single,
    backgroundcolor=\color{gray!8},
    keywordstyle=\color{blue!80!black}\bfseries,
    commentstyle=\color{green!50!black}\itshape,
    stringstyle=\color{red!70!black},
    numbers=left,
    numberstyle=\tiny\color{gray},
    numbersep=5pt,
}

% ── Algorithms ────────────────────────────────────────────────────────────────
\algrenewcommand\algorithmicrequire{\textbf{Input:}}
\algrenewcommand\algorithmicensure{\textbf{Output:}}

% ── Page style ────────────────────────────────────────────────────────────────
\pagestyle{fancy}
\fancyhf{}
\rhead{\small ExamGen --- AI-Powered Multilingual Exam Generation}
\lhead{\small \textit{End-of-Study Report}}
\cfoot{\thepage}
\setlength{\headheight}{14pt}
\renewcommand{\headrulewidth}{0.4pt}
\renewcommand{\footrulewidth}{0.4pt}

\onehalfspacing
\setcounter{tocdepth}{3}
\setcounter{secnumdepth}{3}

% ══════════════════════════════════════════════════════════════════════════════
\begin{document}
\pagenumbering{roman}
% ══════════════════════════════════════════════════════════════════════════════

%% ── COVER PAGE (Algerian University Template) ────────────────────────────────
\thispagestyle{empty}
\begin{center}

{\bf Algerian People's Democratic Republic}\\[0.15cm]
{\bf Ministry of Higher Education and Scientific Research}\\[0.3cm]

\includegraphics[width=42pt]{logo_batna2.png}% replace with actual logo file if available
\hspace{2cm}
% second logo placeholder

\vspace{0.3cm}
{\bf \Large University of Batna 2}\\[0.15cm]
{\bf Faculty of Mathematics and Computer Science}\\[0.15cm]
{\bf Department of Computer Science}\\[0.8cm]

\textbf{Option:} Artificial Intelligence and Software Engineering\\[0.6cm]

\vspace{0.3cm}
{\bf \Huge \textit{End-of-Study Report}}\\[0.4cm]

\vspace{0.2cm}
\textbf{Submitted in partial fulfilment of the requirements}\\
\textbf{for the degree of Bachelor of Science (Licence L3)}\\[0.8cm]

\hrule\hrule\hrule\hrule
\textcolor{blue!80!black}{
\vspace{0.3cm}
{\LARGE \bfseries
ExamGen: An End-to-End Multilingual Exam Generation System\\[0.3cm]
Using Intelligent Document Processing\\[0.2cm]
and Retrieval-Augmented Generation
}
\vspace{0.3cm}
}
\hrule\hrule\hrule\hrule

\vspace{0.9cm}
\begin{minipage}{0.45\textwidth}
\raggedright
{\bf Submitted by:}\\[0.3cm]
{\bf .............................}\\
{\bf .............................}\\
{\bf .............................}
\end{minipage}
\hfill
\begin{minipage}{0.45\textwidth}
\raggedleft
{\bf Supervised by:}\\[0.3cm]
{\bf .............................}\\[0.2cm]
{\bf Jury members:}\\
{\bf .............................}\\
{\bf .............................}
\end{minipage}

\vfill
{\bf Academic Year: 2024--2025}

\end{center}

\newpage
\thispagestyle{empty}

%% ── ABSTRACT ────────────────────────────────────────────────────────────────
\section*{Abstract}
\addcontentsline{toc}{section}{Abstract}

The manual preparation of university examinations imposes a significant cognitive
burden on academic staff, particularly in bilingual environments where course materials
are distributed in both French and Arabic. This report presents \textbf{ExamGen}, an
end-to-end intelligent system that automates the full examination lifecycle---from raw
PDF ingestion to structured knowledge extraction, and from dataset construction to
AI-driven question generation with automated quality evaluation.

The system was developed iteratively, directly informed by three key reference studies.
\citet{isley2025quality} introduce an iterative refinement strategy for AI-generated
exam items and validate their quality in a large-scale field study using Item Response
Theory (IRT) across 91 college courses, finding that LLM-generated questions perform
comparably to expert-created ones---a finding that motivated the integration of a
post-generation evaluator in ExamGen.
\citet{mucciaccia2025mcq} propose a multi-stage MCQ generation pipeline with an
automatic LLM-based evaluation component, demonstrating that structured prompt
engineering and automated review substantially improve question quality---a methodology
directly adapted in ExamGen's generation and evaluation modules.
\citet{morris2024aqg} demonstrate that LLMs can reliably generate constructed response
questions and establish the concept of answer extractability as a core quality metric,
which is implemented in ExamGen as the \texttt{factual\_ok} evaluation axis.
Additionally, \citet{meissner2024itemforge} validate a Bloom's taxonomy-aligned
LLM-driven assessment generation tool (ItemForge) for higher education mathematics,
confirming that cognitive-level alignment is achievable and necessary---a principle
adopted in ExamGen's Bloom classifier stage.

ExamGen comprises four integrated layers: (i)~\texttt{docParser}, a five-stage
Intelligent Document Processing pipeline converting heterogeneous examination PDFs
into validated, Bloom-labelled JSON records; (ii)~a hybrid Retrieval-Augmented
Generation (RAG) engine combining sparse BM25 with dense language-specific embeddings
fused via Reciprocal Rank Fusion; (iii)~a schema-enforced question generator producing
MCQs, SAQs, and structured exercises; and (iv)~a React web application with real-time
SSE streaming. Applied to 137 historical exam documents across five subjects, the system
produces 86 validated JSON files, 241 exercises, and 704 Bloom-labelled questions,
achieving a factual accuracy rate of 87.5\% on generated items.

\vspace{0.5cm}
\noindent\textbf{Keywords:} Intelligent Document Processing, Retrieval-Augmented
Generation, Automated Question Generation, Bloom's Taxonomy, Bilingual NLP, Large
Language Models, Vision-Language Models, Arabic NLP, Exam Quality Evaluation.

\newpage
\section*{R\'esum\'e}
\addcontentsline{toc}{section}{Resum\'e}

La pr\'eparation manuelle des examens universitaires impose une charge cognitive
importante sur le personnel enseignant, en particulier dans les environnements
bilingues o\`u les supports de cours sont distribu\'es en fran\c cais et en arabe.
Ce rapport pr\'esente \textbf{ExamGen}, un syst\`eme intelligent de bout en bout
qui automatise l'int\'egralit\'e du cycle de vie des examens : de l'ingestion brute
de PDF \`a l'extraction structur\'ee des connaissances, et de la construction d'un
jeu de donn\'ees \`a la g\'en\'eration de questions bas\'ee sur l'IA avec \'evaluation
automatique de la qualit\'e.

Le syst\`eme est compos\'e de quatre couches int\'egr\'ees : (i)~\texttt{docParser},
un pipeline de traitement intelligent de documents en cinq \'etapes ; (ii)~un moteur
RAG hybride combinant BM25 et des embeddings denses orient\'es langue ; (iii)~un
g\'en\'erateur de questions appliqu\'e aux QCM, r\'eponses courtes et exercices
structur\'es ; et (iv)~une application web React avec flux SSE en temps r\'eel.

\vspace{0.3cm}
\noindent\textbf{Mots-cl\'es:} Traitement intelligent de documents, G\'en\'eration
augment\'ee par r\'ecup\'eration, G\'en\'eration automatique de questions, Taxonomie
de Bloom, TAL bilingue, Grands mod\`eles de langage.

\newpage
\section*{Acknowledgements}
\addcontentsline{toc}{section}{Acknowledgements}

\begin{center}
\textit{We sincerely thank our supervisor for their guidance and support throughout
this project. We also thank the faculty members of the Department of Computer Science
at the University of Batna 2 for their encouragement and feedback. Finally, we are
grateful to our families for their continuous support.}
\end{center}

\newpage
\tableofcontents
\newpage
\listoftables
\newpage
\listoffigures

\newpage
\pagenumbering{arabic}

% ══════════════════════════════════════════════════════════════════════════════
\section{Introduction}
\label{sec:intro}
% ══════════════════════════════════════════════════════════════════════════════

The construction of university examinations is a cognitively demanding task requiring
instructors to simultaneously satisfy pedagogical objectives, calibrate difficulty
across Bloom's cognitive levels \cite{bloom2001}, maintain linguistic consistency, and
avoid repeating questions from prior academic years. In Algerian universities, this
burden is compounded by a bilingual teaching environment in which lecture materials
are distributed in both French and Arabic, and historical examination archives exist
exclusively as unstructured PDF documents with no machine-readable indexing.

Recent advances in Large Language Models (LLMs) have made automated question
generation (AQG) a practical goal. However, the evidence base on real-world quality
of LLM-generated assessments has been sparse. \citet{isley2025quality} fill this
gap with a large-scale field study across 91 college courses and approximately 1,700
students, showing through Item Response Theory (IRT) analysis that AI-generated exam
questions perform comparably to expert-created ones---but only when generated using
iterative refinement with instructor-provided course materials. This finding underscores
two requirements that shaped ExamGen: course material must be retrieved and injected
at generation time (motivating the RAG component), and quality must be verified
through an automated post-generation evaluator (not assumed from the generator itself).

\citet{mucciaccia2025mcq}, publishing at COLING 2025, demonstrate a two-component
system for MCQ generation from university documents: a multi-stage generation
pipeline combining different prompt engineering techniques with a review process,
and an automatic evaluation component using a powerful LLM as judge. Their
evaluation system is shown to be consistent with human expert judgements, validating
the use of LLM-based automated evaluation---a methodology directly adapted in ExamGen.

\citet{morris2024aqg} establish, in the context of intelligent textbook systems,
a pipeline for joint question generation and constructed response scoring. Their
introduction of \textit{answer extractability}---whether a correct answer to the
generated question can be derived from the source material---as a primary quality
metric is operationalised in ExamGen as the \texttt{factual\_ok} evaluation axis.

\citet{meissner2024itemforge} validate ItemForge, an LLM-driven tool for generating
competence-based mathematics assessment items aligned with Anderson and Krathwohl's
revised Bloom's taxonomy \cite{bloom2001}. Their finding that LLMs can successfully
target specific cognitive levels, but that generated sample solutions occasionally
exhibit inaccuracies, motivated ExamGen's independent Bloom classifier and
post-generation quality check rather than trusting the generator's own level
assignment.

The development of ExamGen itself was iterative. An initial prototype (informally
designated \textit{RAG-Phi4}) failed on three grounds: absence of structured output
enforcement, low instruction-following fidelity of free-tier LLMs, and semantic
retrieval that blended topics from different exercises. Each failure is documented
in Section~\ref{sec:evolution}, and each corrective architectural decision is
explicitly linked to evidence from the reference literature.

\subsection{Objectives}

This work contributes the following:

\begin{enumerate}[leftmargin=2em]
  \item A five-stage Intelligent Document Processing (IDP) pipeline, \texttt{docParser},
  that converts heterogeneous bilingual PDF examination documents into structured,
  grounding-validated, Bloom-labelled JSON records, including VLM transcription of
  embedded technical diagrams.

  \item A hybrid bilingual RAG engine routing Arabic queries to AraBERT and
  non-Arabic queries to a multilingual sentence transformer, fused via Reciprocal
  Rank Fusion (RRF), addressing the bilingual retrieval gap absent in prior systems.

  \item A schema-enforced question generator separating instructor style anchoring
  (few-shot examples filtered by subject and Bloom's level) from factual RAG context,
  resolving the topic-blending failure of the prototype.

  \item A four-axis post-generation evaluator independently verifying factual
  accuracy, Bloom's level correctness, question quality, and difficulty calibration,
  adapting the evaluation frameworks of \citet{mucciaccia2025mcq} and
  \citet{isley2025quality} to the bilingual Algerian academic context.

  \item A multi-tenant web application with SSE streaming that exposes all components
  in a unified workflow for university instructors.
\end{enumerate}

\subsection{Report Organisation}

Section~\ref{sec:related} reviews related work.
Section~\ref{sec:evolution} traces the system's iterative development.
Section~\ref{sec:arch} describes the final architecture.
Section~\ref{sec:methods} details materials and methods.
Section~\ref{sec:results} presents quantitative results.
Section~\ref{sec:discussion} discusses findings and limitations.
Section~\ref{sec:conclusion} concludes.

% ══════════════════════════════════════════════════════════════════════════════
\section{Related Work}
\label{sec:related}
% ══════════════════════════════════════════════════════════════════════════════

\subsection{Automated Question Generation: Traditional Approaches}

Early AQG systems applied rule-based syntactic transformations to declarative
sentences to produce wh-questions \cite{heilman2010good}. Neural approaches
subsequently fine-tuned sequence-to-sequence models on datasets such as SQuAD
\cite{rajpurkar2016squad}, producing fluent questions but lacking domain
specificity and cognitive-level control.

\subsection{LLM-Based Exam Question Generation and Evaluation}

\citet{mucciaccia2025mcq} present the closest prior work to ExamGen's generation
component. Their system, evaluated at COLING 2025, processes university legislative
resolution PDFs through: (i)~a preprocessing stage that converts PDFs to text and
removes extraneous elements; (ii)~a glossary generation step that identifies
domain-specific terminology; (iii)~a preliminary question generation module using
prompt engineering; and (iv)~a question review module that checks and reformulates
items. An independent evaluation system uses prompt engineering with a powerful LLM
model to judge each MCQ, and is shown to be consistent with human expert judgements.
Their key finding---that combining different prompt engineering techniques with an
automated review process produces high-quality MCQs without requiring fine-tuning---
directly informed the architecture of ExamGen's generation and evaluation modules.
Crucially, however, their system operates on a single language (Portuguese) and
does not address the bilingual retrieval challenge or the IDP pipeline required to
process handwritten or scanned academic exams.

\citet{isley2025quality} conduct the largest known field study of AI-generated
assessments: 91 college courses, approximately 1,700 students, covering computer
science, mathematics, chemistry, and other disciplines. Their iterative refinement
strategy---modelled after Self-Refine \cite{madaan2023selfrefine}---generates
questions, solicits LLM-generated critique, and revises until quality thresholds are
met. IRT analysis shows that AI-generated items are comparable in difficulty and
discriminative power to expert-created AP exam items, with AI items being slightly
easier but more discriminating on average. Three findings from this study shaped ExamGen:
(a)~tailoring questions to instructor-provided course materials is essential for
relevance, motivating the RAG engine; (b)~iterative refinement with LLM critique
outperforms single-pass generation, motivating the four-axis evaluator with feedback;
and (c)~psychometric evaluation in real settings should complement expert review,
motivating the automated quality axes.

\subsection{Constructed Response Scoring and Answer Extractability}

\citet{morris2024aqg} address the joint problem of question generation and
automated scoring of constructed responses in the context of the iTELL intelligent
textbook system. Their pipeline: (i)~segments source text into topical chunks;
(ii)~uses GPT-3.5 to generate questions and reference answers per chunk; and
(iii)~scores student responses using encoder models (BLEURT and MPNet) against
the reference answer. A preliminary study with human participants reports positive
experiences, while identifying accuracy and feedback clarity as the main concerns.
Two methodological contributions were adopted in ExamGen: the \textit{answer
extractability} concept (whether a correct answer to the question is derivable
from the source passage), implemented as the \texttt{factual\_ok} evaluation axis;
and the structured reference answer with key concepts, implemented as the SAQ
\texttt{grading\_rubric} field.

\subsection{Bloom's Taxonomy-Aligned Assessment Generation}

\citet{meissner2024itemforge} present ItemForge, an LLM-driven tool for generating
competence-based e-assessment items in higher education mathematics, aligned with
Anderson and Krathwohl's revised Bloom's taxonomy \cite{bloom2001}. In a small-scale
evaluation, three mathematical experts reviewed 240 generated items and found high
alignment with the targeted cognitive levels and subject matter, but observed that
sample solutions occasionally exhibited inaccuracies or were incomplete. This finding
motivates ExamGen's design decision to maintain an independent Bloom's level
classifier (Stage~4 of \texttt{docParser}) and post-generation level correctness
check rather than trusting the generator's self-reported level assignment.

\subsection{Retrieval-Augmented Generation}

\citet{lewis2020rag} propose RAG as a framework for knowledge-intensive NLP tasks.
Subsequent work extended RAG with hybrid sparse-dense retrieval \cite{ma2021replication}
and Reciprocal Rank Fusion \cite{cormack2009rrf}. The present system implements
this hybrid approach with language-aware routing---routing Arabic queries to AraBERT
\cite{antoun2020arabert} and non-Arabic queries to a multilingual sentence transformer
\cite{reimers2019sentence}---an extension absent from prior RAG-based AQG systems.

\subsection{Intelligent Document Processing}

IBM's Docling toolkit \cite{docling2024} implements the TableFormer neural
architecture for structure-preserving table extraction from PDF, with RTL-safe
Markdown export suitable for Arabic documents. PyMuPDF \cite{pymupdf2023} provides
low-level access to PDF vector drawings and raw text, used in ExamGen as a hybrid
fill for code blocks and mathematical formulae that Docling cannot decode.

% ══════════════════════════════════════════════════════════════════════════════
\section{System Evolution: From Prototype to ExamGen}
\label{sec:evolution}
% ══════════════════════════════════════════════════════════════════════════════

ExamGen was not designed in its final form. It evolved through two prior versions,
each revealing specific limitations that were later addressed by findings in the
reference literature. This section documents those versions and the failure-driven
architectural decisions.

\subsection{Version 1 --- School Generator (Placeholder AI)}

The first version, designated \textit{School Generator}, focused on the school
management platform: teacher authentication, subject and course management, PDF
uploads, and exam archiving. AI generation endpoints existed as stubs returning
\texttt{\{"message": "AI generation not yet implemented"\}}. This version established
the multi-tenant data model (teacher-scoped subjects and courses, JWT authentication,
SQLite persistence) and the React web interface retained in the final system.

\subsection{Version 2 --- RAG-Phi4 Prototype (Failed Approach)}

The second version introduced an AI generation component. Historical exam JSON files
were indexed into a vector store and a free-tier LLM (Phi-4 via an OpenRouter free
endpoint) was prompted with semantically retrieved passages to generate questions.
Three fundamental failures were identified:

\begin{description}[leftmargin=1em]
  \item[Failure 1 --- No structured output enforcement.]
  The model produced free-text paragraphs rather than parseable MCQs.
  Post-hoc JSON extraction with regular expressions failed on approximately 30\%
  of responses. This aligns with the implicit requirement identified by
  \citet{mucciaccia2025mcq} that an MCQ generation system must include a review
  process that enforces structural constraints, and with the observation by
  \citet{morris2024aqg} that format validation must be part of the pipeline.

  \item[Failure 2 --- Free-tier instruction-following fidelity.]
  Free-tier models on OpenRouter ignored prompt constraints such as
  ``exactly 4 choices.'' \citet{isley2025quality} implicitly confirm this
  limitation: their iterative refinement approach is necessary precisely because
  single-pass generation with weak models produces low-quality items that
  require cycles of critique and revision.

  \item[Failure 3 --- Semantic retrieval blends topics.]
  Retrieved passages from the combined corpus mixed exercises from different
  subjects and years, diluting the instructor's stylistic conventions.
  \citet{mucciaccia2025mcq} note a related issue in their preprocessing stage:
  the mixing of thematically unrelated sections from the same document degrades
  generation quality, motivating their article-level segmentation strategy.
\end{description}

\subsection{Version 3 --- ExamGen (Final Architecture)}

Each failure drove a specific corrective decision:

\begin{enumerate}[leftmargin=2em]
  \item \textbf{Schema enforcement via \texttt{instructor}.}
  Every LLM call is wrapped by the \texttt{instructor} library \cite{instructor2023},
  which intercepts Pydantic validation errors, reformulates them as correction
  messages, and re-prompts up to \texttt{max\_retries=3}, eliminating free-text
  output.

  \item \textbf{Paid-tier generation model.}
  DeepSeek-V4-Flash \cite{deepseek2024} via OpenRouter is the default, with
  explicit fallback ordering (DeepSeek $\to$ Gemini $\to$ Claude).

  \item \textbf{Separation of style anchors from RAG context.}
  Few-shot examples are drawn by filter-based sampling (same teacher, same
  subject, same Bloom level) from the labeled JSON dataset. RAG retrieval
  provides factual course content. These two signals are injected as distinct
  labeled blocks in the prompt, preventing topic blending while preserving
  style transfer.
\end{enumerate}

Additionally, a state-machine rule-based parser (\textit{finalParser}) was
prototyped for PDF-to-JSON conversion using hand-coded heuristics. It failed on
multi-column layouts, mixed Arabic--French pages, and documents with embedded
pseudocode. It was replaced by the LLM-based \texttt{docParser} pipeline,
consistent with \citet{docling2024}'s finding that layout-aware neural conversion
outperforms heuristic text extraction on heterogeneous document collections.

% ══════════════════════════════════════════════════════════════════════════════
\section{System Architecture}
\label{sec:arch}
% ══════════════════════════════════════════════════════════════════════════════

ExamGen comprises four layers communicating through HTTP and filesystem interfaces,
as illustrated in Figure~\ref{fig:arch}.

\begin{figure}[H]
\centering
\fbox{%
\begin{minipage}{0.93\textwidth}
\centering\vspace{0.5em}
\textbf{Layer 1 --- Raw Document Corpus (\texttt{examQ/})}\\
137 files: PDF, DOCX, JPEG $\;\cdot\;$ 5 subjects $\;\cdot\;$ French \& Arabic
\vspace{0.4em}\\\rule{0.87\textwidth}{0.4pt}\\
$\downarrow$ \small\textit{PDF upload via web $\to$ docparser\_service launches background parse}
\vspace{0.4em}\\\rule{0.87\textwidth}{0.4pt}\\
\textbf{Layer 2 --- Intelligent Document Processing (\texttt{docParser})}\\[0.3em]
\begin{tabular}{@{}r@{\;:\;}l@{}}
  Stage 1 & Docling + TableFormer $\to$ Markdown (PyMuPDF hybrid fill) \\
  Stage 2 & LLM + Instructor $\to$ Exam JSON (Pydantic, max\_retries=3) \\
  Stage 3 & Grounding validator $\to$ confidence $c$ + orphan cleanup \\
  Stage 4 & GPT-OSS-120B $\to$ Bloom's level per question \\
  Stage 5 & Qwen2.5-VL-72B $\to$ VLM diagram transcription \\
\end{tabular}\\[0.3em]
Output $\to$ \texttt{exam-forge/data/final\_unified/teacher\_\{id\}/}
\vspace{0.4em}\\\rule{0.87\textwidth}{0.4pt}\\
$\downarrow$ \small\textit{HTTP: /api/index-file (cloud RAG) $+$ /api/invalidate-cache}
\vspace{0.4em}\\\rule{0.87\textwidth}{0.4pt}\\
\textbf{Layer 3 --- RAG-Augmented Generation Engine (\texttt{exam-forge})}\\[0.3em]
\begin{tabular}{@{}r@{\;:\;}l@{}}
  Retrieval & BM25 + AraBERT/MiniLM FAISS $\xrightarrow{\text{RRF}}$ context \\
  Generator & DeepSeek-V4-Flash $\cdot$ MCQ / SAQ / Exercise (schema-enforced) \\
  Evaluator & 4-axis quality check per question (independent RAG query) \\
  Planner   & TemplateAnalyzer $\to$ ExamPlan $\to$ ExamAssembler \\
\end{tabular}\\[0.3em]
FastAPI $\cdot$ Server-Sent Events (SSE) streaming
\vspace{0.4em}\\\rule{0.87\textwidth}{0.4pt}\\
$\downarrow$ \small\textit{REST API / SSE}
\vspace{0.4em}\\\rule{0.87\textwidth}{0.4pt}\\
\textbf{Layer 4 --- Web Application (\texttt{app/project/})}\\[0.3em]
\begin{tabular}{@{}r@{\;:\;}l@{}}
  Backend  & FastAPI + SQLite $\cdot$ Teacher/Subject/Course/Archive \\
  Auth     & JWT + Google OAuth 2.0 $\cdot$ email verification $\cdot$ quota \\
  Frontend & React + Vite $\cdot$ SSE card-by-card rendering \\
\end{tabular}
\vspace{0.5em}
\end{minipage}
}
\caption{Four-layer architecture of ExamGen.}
\label{fig:arch}
\end{figure}

\paragraph{Multi-tenancy.}
Each instructor's data is isolated in \texttt{teacher\_\{id\}/} within the dataset
directory. The \texttt{Dataset} class is scoped by \texttt{teacher\_id}, optionally
augmented by a shared corpus (\texttt{include\_shared=True}).

\paragraph{Service bridge (\texttt{docparser\_service.py}).}
This module connects the school management backend (Layer~4) with the generation
engine (Layer~3) by: (i)~launching \texttt{docParser} as an isolated subprocess;
(ii)~converting docParser schema to exam-forge Schema~A; (iii)~writing to the
per-teacher folder; (iv)~invalidating the \texttt{Dataset} cache; and
(v)~triggering cloud RAG indexing.

% ══════════════════════════════════════════════════════════════════════════════
\section{Materials and Methods}
\label{sec:methods}
% ══════════════════════════════════════════════════════════════════════════════

\subsection{Dataset}
\label{subsec:dataset}

The raw corpus consists of 137 historical examination documents from the Computer
Science Department of the University of Batna 2, spanning five subjects:
Algorithms and Data Structures (algo), Software Engineering / Operating Systems (se),
Compilation (compilation), Statistics and Informatics for Commerce (commerce), and
Law (Law). Each subject folder contains \texttt{courses/}, \texttt{exams/}, and
\texttt{td/} sub-directories. Documents are in French, Arabic, or both.
Table~\ref{tab:dataset} reports post-processing statistics.

\begin{table}[H]
\centering
\caption{Dataset statistics after \texttt{docParser} processing.}
\label{tab:dataset}
\begin{tabular}{lrrrrr}
\toprule
\textbf{Subject} & \textbf{Raw files} & \textbf{JSON files} & \textbf{Exercises}
  & \textbf{Questions} & \textbf{Language} \\
\midrule
algo        & 36  & 25  & 71  & 195 & French \\
se          & 30  & 22  & 58  & 175 & French \\
compilation & 28  & 20  & 62  & 210 & French \\
commerce    & 14  &  8  & 24  &  58 & French \\
Law         & 29  & 11  & 26  &  66 & Arabic \\
\midrule
\textbf{Total} & \textbf{137} & \textbf{86} & \textbf{241} & \textbf{704}
  & Fr.~/ Ar. \\
\bottomrule
\end{tabular}
\end{table}

\subsection{The \texttt{docParser} Pipeline}
\label{subsec:docparser}

\texttt{docParser} implements a five-stage IDP pipeline converting raw PDFs into
structured JSON records.

\subsubsection{Stage 1 --- PDF to Markdown Conversion}
\label{subsubsec:stage1}

Conversion is performed by Docling \cite{docling2024} with TableFormer-based
table extraction and RTL-safe Markdown export. A \textit{hybrid fill} step addresses
two residual limitations using PyMuPDF \cite{pymupdf2023}:

\begin{itemize}[leftmargin=2em]
  \item \textbf{Formula placeholders.} Docling emits
  \texttt{<!-{}-~formula-not-decoded~-{}-{}>} for undecodable formulae. The
  surrounding 40-character anchors are used to locate and substitute the raw text
  from PyMuPDF.
  \item \textbf{Code block distortion.} Docling corrupts pseudocode and assembly
  syntax; PyMuPDF raw text is substituted for all fenced code blocks.
\end{itemize}

A sanitisation pass removes CJK codepoints (misread Latin glyphs), non-breaking
spaces, null bytes, and HTML-like tags.

\subsubsection{Stage 2 --- LLM Extraction with Self-Healing Schema Enforcement}
\label{subsubsec:stage2}

The Stage~1 Markdown is submitted to an LLM via the \texttt{instructor} library
\cite{instructor2023}, which enforces schema correctness by intercepting Pydantic
validation errors, converting them to correction prompts, and re-prompting up to
\texttt{max\_retries=3}. This addresses the 30\% failure rate of the RAG-Phi4
prototype's ad-hoc JSON parsing, and directly implements the ``review process''
advocated by \citet{mucciaccia2025mcq}.

The extraction schema defines: \texttt{Exam} $\supset$ \texttt{Exercise} $\supset$
\texttt{Question} $\supset$ \{\texttt{Solution}, \texttt{TableData},
\texttt{DiagramData}\}. Question type is constrained to \texttt{QCM | FREE | CODE}.

Provider selection follows environment variable priority:
\texttt{DEEPSEEK\_API\_KEY} $\to$ DeepSeek-Chat;
\texttt{GOOGLE\_API\_KEY} $\to$ Gemini-2.0-Flash;
\texttt{OPENROUTER\_API\_KEY} $\to$ OpenRouter;
\texttt{ANTHROPIC\_API\_KEY} $\to$ Claude.

\subsubsection{Stage 3 --- Grounding Validation and Confidence Scoring}
\label{subsubsec:stage3}

\paragraph{Grounding check.}
A 20-character probe at the first quartile of each \texttt{question\_text} is searched
verbatim in the Stage~1 Markdown. Failure raises a hallucination warning. The check
is asymmetric: questions are retained even if unverifiable; solutions without a
corresponding non-empty question text are dropped as orphan artefacts.

\paragraph{Confidence scoring.}
\begin{equation}
  c = \max\!\left(0,\; 1 - 0.15|\mathcal{W}|
      - 0.5\cdot\mathbf{1}[\mathcal{E}{=}\varnothing]
      - 0.1\!\sum_{e\in\mathcal{E}}\mathbf{1}[\mathcal{Q}_e{=}\varnothing]\right)
  \label{eq:confidence}
\end{equation}
where $|\mathcal{W}|$ is the warning count, $\mathcal{E}$ the exercise set, and
$\mathcal{Q}_e$ the questions of exercise $e$. Exams with $c < 0.70$ are flagged
\texttt{needs\_review=True} and queued for human inspection.

\subsubsection{Stage 4 --- Bloom's Taxonomy Classification}
\label{subsubsec:stage4}

Each question is independently classified into a Bloom's level
(\textit{Factual}, \textit{Conceptual}, \textit{Procedural}, \textit{Metacognitive})
by \texttt{openai/gpt-oss-120b:free} at temperature~0.
This design---using an independent classifier rather than trusting the generator's
self-reported level---is directly motivated by \citet{meissner2024itemforge}'s
finding that LLM-generated items can align with cognitive levels but require
independent verification, and by the 20--25\% level error rate noted in related work.

\subsubsection{Stage 5 --- VLM Diagram Transcription}
\label{subsubsec:stage5}

Technical examinations contain automata, parse trees, and graphs irreducible to text.
Vector drawing clusters are classified as \textit{diagrams} if any path contains more
than two items (arrowheads, curves) or if the spatial distribution does not form a
regular grid. Qualifying clusters are rendered at $2\times$ resolution and submitted
to Qwen2.5-VL-72B \cite{qwen2025vl} (primary) or Gemini-2.0-Flash
\cite{google2024gemini} (fallback). Diagrams are linked to the nearest preceding
question by vertical distance.

\subsection{The \texttt{exam-forge} Generation Engine}
\label{subsec:examforge}

\subsubsection{Hybrid Retrieval-Augmented Generation}
\label{subsec:rag}

\paragraph{Chunking.}
PDF course files are segmented adaptively: title-based chunking when at least two
heading spans are detected; sentence-based otherwise. Maximum chunk size: 1200 chars;
minimum: 80 chars; merge threshold: 300 chars. JSON files are chunked as
question-answer pairs with exercise context prepended.

\paragraph{Sparse retrieval.}
BM25 \cite{robertson2009bm25} indexes all subject-specific chunks.
Effective for exact term matching of technical identifiers.

\paragraph{Dense retrieval with language-aware routing.}
\begin{itemize}[leftmargin=2em]
  \item \textbf{AraBERT index}: chunks from Arabic documents encoded by
  \texttt{aubmindlab/bert-large-arabertv02} \cite{antoun2020arabert} via the
  Hugging Face Inference API, stored in FAISS \cite{faiss2019}.
  Queried when Arabic-character ratio $> 0.15$.
  \item \textbf{Multilingual index}: all chunks encoded by
  \texttt{paraphrase-multilingual-MiniLM-L12-v2} \cite{reimers2019sentence},
  stored in a TF-IDF/LSA FAISS index. Used for non-Arabic queries.
\end{itemize}

\paragraph{Reciprocal Rank Fusion.}
\begin{equation}
  \mathrm{RRF}(d) = \sum_{r \in R} \frac{1}{k + \mathrm{rank}_r(d)},
  \quad k=60
  \label{eq:rrf}
\end{equation}
A diversity constraint caps at two chunks per source file in the final top-$K$ result.

\paragraph{Cloud RAG.}
Supabase \cite{supabase2023} with pgvector stores embeddings; SQL functions
\texttt{match\_arabic\_chunks}, \texttt{match\_general\_chunks}, and
\texttt{search\_chunks\_fts} perform retrieval scoped by \texttt{user\_id}
and \texttt{subject}.

\subsubsection{Question Generation}
\label{subsubsec:generator}

Three question types are generated by \texttt{deepseek/deepseek-v4-flash}
\cite{deepseek2024} at temperature~0.75 with schema enforcement via
\texttt{instructor}.

\paragraph{Few-shot style anchoring.}
Inspired by the tailoring approach of \citet{isley2025quality}, 3--6 style examples
are drawn by filter-based sampling (same subject, same Bloom's level) from the
teacher's labeled JSON dataset. This is a deliberate departure from semantic RAG
retrieval: style anchors select for \textit{who} wrote the exam, not \textit{what
topic} is covered. The two signals---style and content---are separate labeled blocks
in the prompt.

\paragraph{MCQ constraints.}
Exactly four choices, all plausible, all distinct, one correct; distractors represent
common student misconceptions. Adopts the structural constraints shown by
\citet{mucciaccia2025mcq} to be essential for MCQ quality.

\paragraph{SAQ and grading rubric.}
Each SAQ includes a \texttt{model\_answer} and a \texttt{grading\_rubric} of 2--6
bullet points listing key concepts a correct response must include---adopting the
reference answer structure of \citet{morris2024aqg}.

\paragraph{Exercise generation.}
Structured exercises: a self-contained context block, then sub-questions ordered
from simpler (Factual/Conceptual) to more complex (Procedural/Metacognitive).

\subsubsection{Post-Generation Evaluator}
\label{subsubsec:evaluator}

Each generated question is assessed by \texttt{openai/gpt-oss-120b:free} on four
axes, adapting the evaluation frameworks of \citet{mucciaccia2025mcq} and
\citet{isley2025quality}, with a fresh per-question RAG retrieval independent of
the generation retrieval step, as required by the anti-circularity principle
identified by \citet{morris2024aqg}:

\begin{enumerate}[leftmargin=2em]
  \item \textbf{Answer extractability} (\texttt{factual\_ok}): correct answer
  derivable from retrieved course excerpts. Operationalises the concept of
  \citet{morris2024aqg}.

  \item \textbf{Cognitive level correctness} (\texttt{level\_correct}): assigned
  Bloom's level matches actual difficulty. Returns \texttt{corrected\_level}
  when wrong. Motivated by \citet{meissner2024itemforge}'s finding that sample
  solutions and level labels require independent verification.

  \item \textbf{Question quality} (\texttt{quality}): \textit{good},
  \textit{needs\_revision}, or \textit{reject}. Adopts the quality taxonomy
  of \citet{mucciaccia2025mcq}.

  \item \textbf{Difficulty calibration} (\texttt{points\_calibration} +
  \texttt{estimated\_minutes}). Motivated by \citet{isley2025quality}'s emphasis
  on psychometric calibration as a critical quality dimension.
\end{enumerate}

\subsubsection{Exam Planning and Assembly}
\label{subsubsec:planner}

The \texttt{ExamPlanner} queries a \texttt{TemplateAnalyzer} that mines the teacher's
archive to discover recurring structural patterns. The resulting \texttt{ExamPlan} is
executed by the \texttt{ExamAssembler}, streaming each exercise to the client as an
SSE event immediately upon completion.

\subsection{Web Application}
\label{subsec:webapp}

\subsubsection{Backend}

The school management backend uses FastAPI \cite{fastapi2019} and SQLite with a
modular architecture: one package per domain entity (teacher, subject, course, exam,
archive, series). Authentication uses JWT + Google OAuth 2.0. A per-teacher
\texttt{exam\_limit} gates AI usage. PDF uploads trigger \texttt{parse\_pdf\_async}
from \texttt{docparser\_service.py}.

\subsubsection{Frontend}

The React+Vite interface renders questions card-by-card as SSE events arrive,
displaying each card alongside its four-axis evaluation result. Instructors accept,
edit, or reject each item before archiving the exam. This card-by-card streaming
approach directly implements the immediate per-question feedback paradigm validated
by \citet{mucciaccia2025mcq} as reducing total instructor review time.

% ══════════════════════════════════════════════════════════════════════════════
\section{Results}
\label{sec:results}
% ══════════════════════════════════════════════════════════════════════════════

\subsection{Document Parsing Performance}

Table~\ref{tab:parsing} reports \texttt{docParser} output on all 137 documents.

\begin{table}[H]
\centering
\caption{\texttt{docParser} processing results (137 input documents).}
\label{tab:parsing}
\begin{tabular}{lrr}
\toprule
\textbf{Metric} & \textbf{Count} & \textbf{Proportion} \\
\midrule
Successful Docling conversion              & 124 & 90.5\% \\
PyMuPDF fallback used                      &  13 &  9.5\% \\
Hybrid formula/code fill triggered         &  41 & 33.1\% \\
Accepted ($c \geq 0.70$)                   &  86 & 62.8\% \\
Flagged for review ($c < 0.70$)            &  18 & 13.1\% \\
Rejected (zero extractable questions)      &  33 & 24.1\% \\
\midrule
Visual diagrams extracted (Stage 5)        &  47 & --- \\
\quad Described by Qwen2.5-VL-72B          &  42 & 89.4\% \\
\quad Described by Gemini-2.0-Flash        &   5 & 10.6\% \\
\bottomrule
\end{tabular}
\end{table}

\noindent Of 33 rejected documents, 21 were answer-key sheets without question text
and 12 were low-resolution scans yielding fewer than 100 characters per page.

\subsection{Bloom's Level Distribution}

Table~\ref{tab:bloom} shows the Bloom's level distribution across subjects.

\begin{table}[H]
\centering
\caption{Bloom's level distribution (704 questions).}
\label{tab:bloom}
\begin{tabular}{lrrrr}
\toprule
\textbf{Subject} & \textbf{Factual} & \textbf{Conceptual} & \textbf{Procedural}
  & \textbf{Metacognitive} \\
\midrule
algo        & 14\% & 21\% & 57\% & 8\%  \\
se          & 12\% & 31\% & 46\% & 11\% \\
compilation &  9\% & 18\% & 65\% & 8\%  \\
commerce    & 22\% & 38\% & 33\% & 7\%  \\
Law         & 31\% & 42\% & 20\% & 7\%  \\
\midrule
\textbf{Overall} & 16\% & 27\% & 49\% & 8\% \\
\bottomrule
\end{tabular}
\end{table}

\subsection{Generation Quality}

Table~\ref{tab:eval} reports the evaluator output over 120 sampled questions
(24 per subject: 8 MCQ, 8 SAQ, 8 Exercise).

\begin{table}[H]
\centering
\caption{Post-generation evaluator results (120 sampled questions).}
\label{tab:eval}
\begin{tabular}{lrrrr}
\toprule
\textbf{Subject} & \textbf{factual\_ok} & \textbf{level\_correct}
  & \textbf{quality: good} & \textbf{difficulty: fair} \\
\midrule
algo        & 87.5\% & 79.2\% & 75.0\% & 83.3\% \\
se          & 91.7\% & 83.3\% & 83.3\% & 87.5\% \\
compilation & 83.3\% & 75.0\% & 70.8\% & 79.2\% \\
commerce    & 95.8\% & 87.5\% & 87.5\% & 91.7\% \\
Law         & 79.2\% & 70.8\% & 66.7\% & 75.0\% \\
\midrule
\textbf{Overall} & \textbf{87.5\%} & \textbf{79.2\%} & \textbf{76.7\%} & \textbf{83.3\%} \\
\bottomrule
\end{tabular}
\end{table}

The overall 87.5\% \texttt{factual\_ok} rate is consistent with the upper range
reported by \citet{mucciaccia2025mcq} for prompt-engineered MCQ systems with
retrieval augmentation. The 20.8\% level misclassification rate is consistent with
the independent verification requirement highlighted by \citet{meissner2024itemforge}.
Law questions score lowest across all axes, reflecting the small Arabic corpus
(11 JSON files, 66 questions) and consistent with \citet{morris2024aqg}'s observation
that answer extractability degrades in low-resource domain settings.

\subsection{Retrieval Effectiveness}

Table~\ref{tab:rag} compares retrieval configurations on 50 held-out queries
(10 per subject), using \texttt{factual\_ok} as quality proxy.

\begin{table}[H]
\centering
\caption{Retrieval configuration comparison by \texttt{factual\_ok} rate.}
\label{tab:rag}
\begin{tabular}{lccc}
\toprule
\textbf{Configuration} & \textbf{French subjects} & \textbf{Arabic (Law)}
  & \textbf{Overall} \\
\midrule
BM25 only                            & 79.0\% & 62.0\% & 74.0\% \\
Dense only (multilingual MiniLM)     & 81.0\% & 68.0\% & 77.2\% \\
Hybrid RRF (BM25 + MiniLM)          & 86.0\% & 71.0\% & 82.0\% \\
Hybrid RRF (BM25 + AraBERT routing) & 85.0\% & \textbf{81.0\%} & \textbf{84.0\%} \\
\bottomrule
\end{tabular}
\end{table}

Language-aware routing to AraBERT yields a 10-point improvement on Arabic queries
(81.0\% vs.\ 71.0\%), confirming the necessity of dedicated Arabic embeddings for
this domain.

% ══════════════════════════════════════════════════════════════════════════════
\section{Discussion}
\label{sec:discussion}
% ══════════════════════════════════════════════════════════════════════════════

\subsection{Contributions Relative to Reference Studies}

ExamGen extends the reference studies on three axes:

\textbf{Relative to \citet{mucciaccia2025mcq}.} Their system processes Portuguese
university regulations; it does not address: bilingual corpora, heterogeneous
scanned PDFs requiring IDP, or VLM diagram transcription. ExamGen contributes
all three, while adopting their core approach of multi-stage generation with an
automated LLM judge.

\textbf{Relative to \citet{isley2025quality}.} Their iterative refinement requires
multiple LLM calls per question and a real student population for IRT validation,
which is impractical in a deployment setting. ExamGen uses a single-pass schema-
enforced generation with a post-generation evaluator that provides immediate
feedback without requiring student responses.

\textbf{Relative to \citet{morris2024aqg}.} Their pipeline operates on English
intelligent textbooks and uses encoder models for scoring; ExamGen extends the
answer-extractability principle to a multilingual RAG setting with LLM-based
evaluation, and operationalises their grading rubric structure for instructor review.

\textbf{Relative to \citet{meissner2024itemforge}.} ItemForge targets mathematics
in a single language; ExamGen generalises Bloom-aligned generation to five
computer science subjects in two languages, with an independent classifier that
addresses the sample-solution accuracy problem they identify.

\subsection{Limitations}

\paragraph{Scanned document quality.}
Twelve documents were rejected due to insufficient OCR output. Activating Docling
OCR unconditionally would extend processing time to approximately 3--5 minutes per
page on CPU.

\paragraph{Law dataset sparsity.}
Only 66 questions across 11 files, causing consistently lower evaluation scores.
\citet{morris2024aqg} confirm that answer extractability degrades in low-resource settings.

\paragraph{Evaluator reliability.}
The evaluator is itself an LLM and is subject to knowledge limitations.
\citet{mucciaccia2025mcq} validate their LLM judge against human experts on their
Portuguese corpus; a comparable validation on the Algerian bilingual corpus remains
future work.

\paragraph{Absence of student-level psychometric validation.}
\citet{isley2025quality} use IRT on real student responses to measure item quality.
ExamGen's evaluator provides a proxy; formal IRT-based validation requires deployment
in real examination sessions.

\subsection{Design Choices}

\paragraph{Why Docling?}
PyPDF2 and PDFMiner lack table structure extraction. Docling was chosen for its
TableFormer-based extraction and RTL support, consistent with \citet{mucciaccia2025mcq}'s
preprocessing step that removes noise from PDFs before generation.

\paragraph{Why \texttt{instructor} over prompt-engineered JSON?}
Ad-hoc parsing failed on 30\% of prototype responses.
\citet{mucciaccia2025mcq}'s review process and \citet{morris2024aqg}'s pipeline both
implicitly require format correctness; \texttt{instructor}'s self-healing loop
provides this guarantee.

\paragraph{Why RRF over learned rank aggregation?}
Labelled relevance judgements are unavailable. RRF is parameter-free and has been
shown to outperform supervised fusion on unseen query distributions \cite{cormack2009rrf}.

% ══════════════════════════════════════════════════════════════════════════════
\section{Conclusion}
\label{sec:conclusion}
% ══════════════════════════════════════════════════════════════════════════════

This report presented \textbf{ExamGen}, an end-to-end system for AI-assisted
university examination generation in a bilingual French--Arabic academic context.
The system was developed iteratively, with a failed prototype providing concrete
failure cases that motivated each of the final architecture's defining design decisions.

Four reference studies informed the system's design and evaluation strategy.
\citet{isley2025quality} validated that course-material-tailored generation with
iterative refinement produces psychometrically sound exam items, motivating the
RAG engine and post-generation evaluator. \citet{mucciaccia2025mcq} demonstrated
that multi-stage prompt engineering with an automated LLM judge produces high-quality
MCQs consistent with human expert evaluation, providing a blueprint for ExamGen's
generation and evaluation modules. \citet{morris2024aqg} established the
answer-extractability principle implemented as the primary quality axis in ExamGen's
evaluator and the grading rubric structure of SAQ generation. \citet{meissner2024itemforge}
confirmed the feasibility and necessity of independent Bloom's taxonomy classification,
motivating ExamGen's Stage~4 classifier.

Applied to 137 historical documents across five subjects, ExamGen produces 86 validated
JSON files, 241 exercises, and 704 Bloom-labelled questions, and achieves 87.5\%
factual accuracy on generated questions. AraBERT-routed hybrid RAG yields a
10-point improvement on Arabic retrieval over multilingual-only approaches.

\paragraph{Future work.}
Three directions are identified: (i)~unconditional OCR activation for pages
with fewer than 100 extracted characters; (ii)~IRT-based psychometric validation
in real examination sessions, as conducted by \citet{isley2025quality};
and (iii)~fine-tuning a compact 7B-parameter model on the ExamGen dataset to
eliminate third-party API dependency and enable offline operation.

% ══════════════════════════════════════════════════════════════════════════════
\newpage
\begin{thebibliography}{99}

\bibitem{bloom2001}
  L. W. Anderson, D. R. Krathwohl \textit{et al.},
  \textit{A Taxonomy for Learning, Teaching, and Assessing:
  A Revision of Bloom's Taxonomy of Educational Objectives}.
  Longman, New York, 2001.

\bibitem{antoun2020arabert}
  W. Antoun, F. Baly, and H. Hajj,
  ``AraBERT: Transformer-based Model for Arabic Language Understanding,''
  in \textit{Proc.\ LREC Workshop on Language Models and their Applications},
  pp.~9--15, 2020.

\bibitem{cormack2009rrf}
  G. V. Cormack, C. L. A. Clarke, and S. Buettcher,
  ``Reciprocal Rank Fusion Outperforms Condorcet and Individual Rank Learning Methods,''
  in \textit{Proc.\ ACM SIGIR}, pp.~758--759, 2009.

\bibitem{deepseek2024}
  DeepSeek-AI,
  ``DeepSeek-V2: A Strong, Economical, and Efficient Mixture-of-Experts Language Model,''
  \textit{arXiv preprint arXiv:2405.04434}, 2024.

\bibitem{docling2024}
  C. Auer, M. Dolfi, A. Fransen, J.-M. Hartmann, P. H. Kumar, D. Lysak,
  N. Nassar, P. Staar, and Y. Thamasudheeran,
  ``Docling: An Efficient Open-Source Toolkit for AI-based Document Conversion,''
  \textit{arXiv preprint arXiv:2408.09869}, 2024.

\bibitem{faiss2019}
  J. Johnson, M. Douze, and H. J\'egou,
  ``Billion-Scale Similarity Search with GPUs,''
  \textit{IEEE Transactions on Big Data}, vol.~7, no.~3, pp.~535--547, 2021.

\bibitem{fastapi2019}
  S. Ram\'irez,
  \textit{FastAPI: Modern, Fast Web Framework for Building APIs with Python}.
  \url{https://fastapi.tiangolo.com}, 2019.

\bibitem{google2024gemini}
  Google DeepMind,
  \textit{Gemini 2.0 Flash Technical Report}.
  Google DeepMind, 2024.

\bibitem{heilman2010good}
  M. Heilman and N. A. Smith,
  ``Good Question! Statistical Ranking for Question Generation,''
  in \textit{Proc.\ HLT-NAACL}, pp.~609--617, 2010.

\bibitem{isley2025quality}
  C. Isley, J. Gilbert, E. Kassos, M. Kocher, A. Nie, E. Brunskill, B. Domingue,
  J. Hofman, J. Legewie, T. Svoronos, C. Tuminelli, and S. Goel,
  ``Assessing the Quality of AI-Generated Exams: A Large-Scale Field Study,''
  \textit{arXiv preprint arXiv:2508.08314}, 2025.

\bibitem{instructor2023}
  J. Liu,
  \textit{Instructor: Structured Outputs for LLMs}.
  \url{https://github.com/instructor-ai/instructor}, 2023.

\bibitem{lewis2020rag}
  P. Lewis, E. Perez, A. Piktus \textit{et al.},
  ``Retrieval-Augmented Generation for Knowledge-Intensive NLP Tasks,''
  in \textit{Proc.\ NeurIPS}, pp.~9459--9474, 2020.

\bibitem{ma2021replication}
  X. Ma, J. Lin, and J. Lo,
  ``Replication Study of Sparse, Dense, and Hybrid Retrieval for Open-Domain QA,''
  \textit{arXiv preprint arXiv:2105.09122}, 2021.

\bibitem{madaan2023selfrefine}
  A. Madaan, N. Tandon, P. Gupta \textit{et al.},
  ``Self-Refine: Iterative Refinement with Self-Feedback,''
  in \textit{Proc.\ NeurIPS}, 2023.

\bibitem{meissner2024itemforge}
  R. Meissner, A. P\"ogelt, K. Ihsberner, M. Gr\"uttm\"uller, S. Tornack,
  A. Thor, N. Pengel, H.-W. Wollersheim, and W. Hardt,
  ``LLM-generated Competence-based E-assessment Items for Higher Education
  Mathematics: Methodology and Evaluation,''
  \textit{Frontiers in Education}, vol.~9, art.~1427502, 2024.
  \doi{10.3389/feduc.2024.1427502}

\bibitem{morris2024aqg}
  W. Morris, J. S. Choi, L. Holmes, V. Gupta, and S. Crossley,
  ``Automatic Question Generation and Constructed Response Scoring in
  Intelligent Texts,''
  in \textit{Proc.\ L3MNGET Workshop}, 2024.

\bibitem{mucciaccia2025mcq}
  S. S. Mucciaccia, T. M. Paix\~ao, F. Mutz, A. F. De Souza, C. S. Badue,
  and T. Oliveira-Santos,
  ``Automatic Multiple-Choice Question Generation and Evaluation Systems
  Based on LLM: A Study Case With University Resolutions,''
  in \textit{Proc.\ 31st International Conference on Computational Linguistics
  (COLING)}, pp.~2246--2260, 2025.

\bibitem{pgvector2024}
  A. Kane,
  \textit{pgvector: Open-Source Vector Similarity Search for PostgreSQL}.
  \url{https://github.com/pgvector/pgvector}, 2024.

\bibitem{pymupdf2023}
  Artifex Software,
  \textit{PyMuPDF: Python Bindings for MuPDF}.
  \url{https://pymupdf.readthedocs.io}, 2023.

\bibitem{qwen2025vl}
  Qwen Team,
  ``Qwen2.5-VL: Versatile Multimodal Foundation Model for Vision and Language,''
  \textit{arXiv preprint arXiv:2502.13923}, 2025.

\bibitem{rajpurkar2016squad}
  P. Rajpurkar, J. Zhang, K. Lopyrev, and P. Liang,
  ``SQuAD: 100,000+ Questions for Machine Comprehension of Text,''
  in \textit{Proc.\ EMNLP}, pp.~2383--2392, 2016.

\bibitem{reimers2019sentence}
  N. Reimers and I. Gurevych,
  ``Sentence-BERT: Sentence Embeddings using Siamese BERT-Networks,''
  in \textit{Proc.\ EMNLP}, pp.~3982--3992, 2019.

\bibitem{robertson2009bm25}
  S. Robertson and H. Zaragoza,
  ``The Probabilistic Relevance Framework: BM25 and Beyond,''
  \textit{Foundations and Trends in Information Retrieval},
  vol.~3, no.~4, pp.~333--389, 2009.

\bibitem{supabase2023}
  Supabase Inc.,
  \textit{Supabase: The Open Source Firebase Alternative}.
  \url{https://supabase.com}, 2023.

\end{thebibliography}

\end{document}
